\documentclass[12pt]{article}
\usepackage[a4paper, total={6.5in, 9.5in}]{geometry}
\usepackage{graphicx}			
\usepackage{amssymb}
\usepackage{amsmath}			
\usepackage{fancyhdr}
\usepackage{enumerate}       
\usepackage{dirtytalk} 
\usepackage[english]{babel}
\usepackage{amsthm}
\usepackage{xcolor}
\usepackage[shortlabels]{enumitem}
\usepackage{titlesec}
\titleformat{\section}{\normalfont\Large\bfseries}{\thesection}{0pt}{\Large}
\setlength{\parindent}{0pt}
\setlength{\parskip}{0.8em}

\newtheorem{theorem}{Theorem}[section]
\newtheorem{corollary}{Corollary}[theorem]
\newtheorem{lemma}[theorem]{Lemma}
\newtheorem{definition}[theorem]{Definition}


\begin{document}
% \renewcommand{\thesection}{Exercise \arabic{section}}


\setcounter{section}{0}
\renewcommand{\thesection}{}


\parindent = 0pt


      {\bf Analysis 1 HW 5} 
%==================

\section{4.10}
{\bf Theorem 4.19} Let $f$ be a continuous mapping of a compact metric space $X$ into a metric space $Y$. Then $f$ is uniformly continuous on $X$. 

\begin{proof}
      Assume, for contradiction, that $f$ is not uniformly continuous. Then,
      there exists $\epsilon>0$ such 
      that for all $\delta>0$ there exists $p\in X$ such that 
      $f(B_{\delta}(p, X))\not\subset B_{\epsilon}(f(p), Y)$. 
      Let $q\in B_{\delta}(p, X)$ such that 
      $f(q)\in f(B_{\delta}(p, X))\setminus B_{\epsilon}(f(p), Y)$. So, we 
      have $d_{X}(p, q)<\delta$, and $d_{Y}(f(p), f(q))>\epsilon$. Construct 
      sequences $(p_{n})$ and $(q_{n})$ corresponding to the sequence
      $\delta_n =\frac{1}{n}$ such that $d_X(p_n, q_n)<\frac{1}{n}$.
      Thus, $d_{X}(p_{n}, q_{n})\to {0}$, $d_{Y}(f(p_n), f(q_n))>\epsilon$. 

      From Theorem 2.37, $\{ p_{n} \}$ has a limit point 
      $l\in X$. Thus, there exists a subsequence $(p_{n_{k}})$ of $(p_{n})$ 
      which converges to $l$. Because $d_{X}(p_{n}, q_{n})\to 0$, it follows 
      that $(q_{n_{k}})\to l$. By the continuity of $f$ over $X$, 
      $\lim_{ x \to l }f(x)=f(l)$. The sequence criterion tells us 
      that $(p_{n_{k}})\to l$ and $(q_{n_{k}})\to l$ imply that 
      $(f(p_{n_{k}}))\to f(l)$ and $(f(q_{n_{k}}))\to f(l)$. However, 
      this is impossible since 
      $d_{Y}(f(p_{n_{k}}), f(q_{n_{k}}))>\epsilon$. This contradiction 
      forces us to conclude $f$ is uniformly continuous. 
\end{proof}

\pagebreak

\section{4.13}
$E$ is a dense subset of $X$.\\
$f:E\to \mathbb{R}$ is uniformly continuous.

{\bf To prove:} $f$ has a continuous extension from $E$ to $X$.


\begin{proof}
      Because $E$ is dense in $X$, every point of $X$ is either in 
      $E$ or a limit point of $E$. So, all points in $X\setminus E$ are 
      limit points of $E$. Consider $p\in X\setminus E$. There exists a 
      sequence $(p_{n})\to p$ in $E$. Let $\epsilon>0$. Since $f$ is 
      uniformly continuous, there exists $\delta$ such that 
      $d_{X}(p_{n}, p_{m})<\delta$ implies 
      $|f(p_{n})-f(p_{m})|<\epsilon$. Since convergent sequences are 
      Cauchy sequences in metric spaces, there exists $N$ such that 
      $n, m>N$ implies $d_{X}(p_{n}, p_{m})<\delta$. Thus, $n, m>N$ 
      implies $|f(p_{n})-f(p_{m})|<\epsilon$, i.e, $(f(p_{n}))$ is 
      Cauchy. Since $\mathbb{R}$ is complete, $(f(p_{n}))$ must 
      converge to some $l_{p} \in\mathbb{R}$. 

      Consider another sequence $(q_{n})\to p$. It follows that 
      $d_{X}(p_{n}, q_{n})\to 0$. Following the same chain of reasoning 
      as in the previous paragraph leads us to conclude 
      $|f(p_{n})-f(q_{n})|\to 0$. Thus, $f(q_{n})\to l_{p}$. We have 
      shown that for every sequence $(q_{n})\to p$, $f(q_{n})\to l_{p}$. 
      Thus, $\lim_{ x \to p }f(x)=l_{p}$. Now, define 
      $F:X\to \mathbb{R}$ by

      $$
      F(x)=
      \begin{cases}
      f(x) & x\in E \\
      l_{x} & x\not\in E
      \end{cases}
      $$

      We will now show that $F$, as defined, is continuous. 
      Let $\epsilon>0$. Let $s\in X$. If $s\in E$, since $f$ is 
      continuous at $s$, there exists $\delta>0$ such that 
      $f(B_{\delta}(s, E))\subset B_{\frac{\epsilon}{2}}(f(s), 
      \mathbb{R})$. If $s\in X\setminus E$, $\lim_{ x \to s }f(x)=F(s)$ 
      by definition, so there exists $\delta>0$ such that 
      $f(B_{\delta}(s, E)\setminus \{ s \})\subset 
      B_{\frac{\epsilon}{2}}(F(s), \mathbb{R})$. In either case, 
      it follows from our definition of $F$ that 
      $F(B_{\delta}(s, X))\subset \overline{B_{\frac{\epsilon}{2}}(F(s), 
      \mathbb{R})}\subset B_{\epsilon}(F(s), \mathbb{R})$. 
      Thus, $F$ is continuous at $s$.
\end{proof}

\pagebreak

\section{4.15}
$f:X\to Y$ is {\it open} if for all open $V\subset X$, $f(V)$ is open in $Y$. 

{\bf To prove:} Every continuous open mapping of $\mathbb{R}$ into $\mathbb{R}$ is monotonic. 
\begin{proof}
      Let $f:\mathbb{R}\to \mathbb{R}$ be continuous and open. FTSOC, 
      assume $f$ is not monotonic. Then, there exist $x_{1}<x_{2}<x_{3}$ 
      such that $f(x_{1})<f(x_{2})>f(x_{3})$ 
      (the other case $f(x_{1})>f(x_{2})<f(x_{3})$ follows similarly). 
      Since $[x_{1}, x_{3}]$ is compact, $f(x)$ attains a maximum 
      value in $[x_{1}, x_{3}]$, say $m$, courtesy the extreme value 
      theorem. We know that $m\ne f(x_{1})$ and $m\ne f(x_{3})$ 
      since $m\ge f(x_{2})$. Thus, $f(x)$ attains the maximum value 
      $m$ in $(x_{1}, x_{3})$. Consider the image $I=f((x_{1}, x_{3}))$. 
      We know that $m=\sup I$ and $m\in I$. This makes it impossible 
      for $I$ to contain a neighborhood of $m$. Thus, $I$ is not open, 
      which contradicts our hypothesis that $f$ is an open map.
      Therefore, $f$ must be monotonic.
\end{proof}

\pagebreak

\section{2.19}
\begin{enumerate}[(a)]
      \item \textbf{To prove:} If $A$ and $B$ are disjoint closed sets in some metric space $X$, they are separated.
            \begin{proof}
                  Since $A$ and $B$ are closed in $X$, $\overline{A}=A$ and $\overline{B}=B$. We are given that $A\cap B=\emptyset$. It follows that $\overline{A}\cap B=\emptyset$ and $A\cap \overline{B}=\emptyset$. 
            \end{proof}
      \item \textbf{To prove:} If $A$ and $B$ are disjoint open sets in some metric space $X$, they are separated.
            \begin{proof}
                  FTSOC, assume $p\in\overline{A}\cap B$. Since 
                  $A\cap B=\emptyset$, $p\not\in A$. So, $p$ is a limit 
                  point of $A$. It follows that for all $\epsilon>0$, 
                  $B_{\epsilon}(p, X)\cap A\ne \emptyset$, i.e, 
                  $B$ does not contain a neighborhood of $p$. This 
                  contradicts our hypothesis that $B$ is open. Therefore, 
                  we must have $\overline{A}\cap B=\emptyset$. Similarly, 
                  we can show that $A\cap \overline{B}=\emptyset$. 
            \end{proof}
      \item \textbf{To prove:} For $p\in X$, $\delta>0$, define $A=\{ q\ |\ d(p, q)<\delta \}$, $B=\{ q \ |\ d(p, q)>\delta \}$. $A$ and $B$ are separated.
            \begin{proof}
                  $A$ is an open ball. Take any point $q\in B$. Let $\gamma\equiv d(q, p)$ and $r\equiv\frac{{\gamma-\delta}}{2}$. Then, $B_{r}(q, X)\subset B$. Thus, $B$ is also open. Clearly, $A$ and $B$ are disjoint. It follows from (b) that $A$ and $B$ are separated.
            \end{proof}
      \item \textbf{To prove:} Every connected metric space with at least two points is uncountable.
            \begin{proof}
                  Let $X$ be a connected metric space. Let 
                  $x, y\in X$, $x\ne y$. Let $r\equiv d(x, y)$. 
                  Let $\alpha\in(0, r)$. 
                  Define $A\equiv\left\{ q\ |\ d(q, x)< \alpha   \right\}$ 
                  and $B\equiv \left\{  q\ | \ d(q, x)> \alpha \right\}$. 
                  Since $A\sqcup B$ cannot be a separation of $X$, 
                  there must exist $z\in X$ such that $d(x, z)=\alpha$. 
                  Thus, with every $\alpha\in(0, r)$ there must be 
                  associated a unique $z_{\alpha}\in X$. We know that 
                  $(0, r)$ is uncountable. It follows that $X$ is 
                  uncountable.
            \end{proof}
\end{enumerate}

\pagebreak

\section{5.1}
$f:\mathbb{R}\to \mathbb{R}$,\\
$|f(x)-f(y)|\leq (x-y)^{2}\ \ \forall \ x, y \in \mathbb{R}$. 

{\bf To prove:} $f$ is constant.

\begin{proof}
      Let $\Delta x\equiv y-x$. Then, $y=x+\Delta x$. This gives us

      \begin{align*}
      & |f(x+\Delta x)-f(x)|\le \Delta x^{2}\\
      \implies  & \frac{{|f(x+\Delta x)-f(x)|}}{|\Delta x|}\le |\Delta x| \\
      \implies  & \lim_{ \Delta x \to 0 } \left|\frac{{f(x+\Delta x)-f(x)}}{\Delta x}\right|\le \lim_{ \Delta x \to 0 } |\Delta x| \\
      \implies  & |f'(x)|\le 0  \\
      \implies  & f'(x)=0.
      \end{align*}

      Thus, $f'(x)=0$ for all $x\in \mathbb{R}$. FTSOC, assume $f$ is not constant, i.e, 
      there exist $x_{1}$ and $x_{2}$ such that $f(x_{1})\ne f(x_{2})$. WLOG, $x_{1}<x_{2}$. Then, from the mean value theorem, there exists $c$ such that $x_{1}<c<x_{2}$ and $f'(c)=\frac{{f(x_{2})-f(x_{1})}}{x_{2}-x_{1}}\ne 0$. However, $f'(x)=0$ for all $x$. This contradiction implies that $f$ is indeed a constant function.
\end{proof}

\end{document}